% Abstract
\begin{abstract}
\vspace{-1cm}
\onehalfspacing

\noindent This work presents an experimental comparative study of three
implementations of the Quicksort algorithm: the classical recursive version, a hybrid
version in which partitioning stops for subarrays with fewer than $M$ elements, which
are then sorted by Insertion Sort, and a hybrid version using median-of-three pivot
selection. The implementations were developed in C++17 on top of a single
parameterised engine, ensuring that all three variants count comparisons and swaps
identically. The algorithms were evaluated on five test datasets --- random, sorted,
reverse sorted, with many repeated elements, and nearly sorted --- across four input
sizes ranging from one thousand to five hundred thousand elements, measuring
execution time, comparison count and swap count. Empirical calibration indicated
$M \approx 40$ as the optimal value, higher than the range of 5 to 25 recommended by
the classical literature, a difference attributed to modern hardware characteristics.
The hybrid versions reduced execution time by 16\% to 45\% compared to the pure
recursive version. A dedicated experiment forced the algorithm's worst case by
combining a sorted array with an inadequate pivot choice, numerically confirming the
theoretical prediction: the ratio between comparison count and $n^2$ converged to
exactly $0.5000$, a difference of approximately two thousand times relative to the
median-of-three version. The critical analysis revealed that, on random data, the
hybrid version performs more comparisons than the recursive version yet remains
faster, showing that counting elementary operations is no longer a reliable predictor
of execution time on modern hardware.

\vspace*{.75cm}

\noindent \textbf{Keywords:} Quicksort, Insertion Sort, Sorting Algorithms,
Complexity Analysis, Experimental Algorithmics, Median-of-three.

\end{abstract}
